\section{Half-budget identity from a self-dual two-sector reduction}
\label{app:half-budget-self-duality}

The accessible-area reduction leaves one primitive budget identity:
\[
  M_{\rm total}=\frac{E_{\rm eff}}{2}.
\]
A clean route is a two-sector decomposition of the finite-cell quadratic
budget,
\[
  E_{\rm eff}=M_\phi+M_p,
\]
where \(M_\phi\) is the phase sector and \(M_p\) the pressure/compression
sector.  Let
\[
  \gamma=\frac{M_\phi}{M_p}.
\]
Then
\[
  M_\phi=\frac{\gamma}{1+\gamma}E_{\rm eff}.
\]
If the reduced action is self-dual under exchange of the phase and pressure
sectors, then \(\gamma=1\), and hence
\[
  M_\phi=\frac{E_{\rm eff}}{2}.
\]
Combining this with the accessible-area projection
\[
  M_{\rm acc}=M_\phi\frac{(R-a)^2}{R^2}
\]
and the radial optimum
\[
  \Lambda_\phi=\frac{M_{\rm acc}}{(R-a)^2}
\]
gives, for the normalized cell \(R=1\),
\[
  \Lambda_\phi=\frac{E_{\rm eff}}{2}.
\]
Thus the half-budget identity is derived within the self-dual two-sector
finite-cell reduction.  The remaining primitive-equation task is to derive the
phase--pressure exchange symmetry from the interior GP/SST dynamics, rather
than postulating \(\gamma=1\).
